\chapter{Discussion and limitations}\label{ch:discussion}

\section{Answering the research questions}

\subsection{Representation robustness under spelling noise}
\emph{How does classification performance change as spelling swaps become more frequent, and how do the neural models compare with the lexical baseline?}

Across all evaluated corruption levels, every model trained on clean text loses accuracy as adjacent interior letter swaps become more frequent, but the linear baseline proves substantially more resilient than either convolutional network. Between clean text ($0\%$ swaps) and the highest corruption target ($20\%$ swaps), \textbf{Word TF-IDF LR} drops by only $0.91$ percentage points ($90.05\%$ to $89.14\%$), compared with $1.86$ percentage points for \textbf{Subword CNN} ($87.18\%$ to $85.32\%$) and $2.31$ percentage points for \textbf{Word CNN} ($86.79\%$ to $84.48\%$). The baseline leads at every swap target (Figure~\ref{fig:core-robustness}), maintains its lead across both additional random swap draws at $20\%$, and confirms the same ordering on the separate confirmation set of 4,000 articles.

A natural structural explanation for this resilience lies in how sparse linear models and convolutional networks with max pooling aggregate evidence across a document. In Word TF-IDF LR, each retained unigram or bigram contributes additively to the class logits across the entire news item of 25 to 50 words. Even at the $20\%$ swap target for eligible words, where $99.1\%$ of test articles contain at least one typo, only $12.93\%$ of all test tokens are altered, leaving roughly $87\%$ of the words in a typical article intact. When a swap invalidates one word and its adjacent bigrams, the linear classifier still sums evidence over dozens of uncorrupted topical terms elsewhere in the headline and description. In a CNN with global max pooling over the sequence, by contrast, each of the 192 filters forwards only its single highest window activation across the entire document. If a swapped word sits inside the window of 3, 4, or 5 words that produced a filter's global maximum, or if replacing it with \texttt{<unk>} alters the winning window, the pooled feature passed to the classification head changes directly. Consistent with this mechanism, our paired transition analysis (Figure~\ref{fig:paired-changes}) shows that Word TF-IDF LR trained on clean text flips its correctness decision on only $215$ test articles at $20\%$ swaps ($6.1\%$ instability across repeated $20\%$ swap draws), whereas both CNNs flip correctness on nearly $500$ articles ($11.3\%$--$11.8\%$ instability across draws).

Inspecting errors at the topic and article level reveals that aggregate accuracy drops tell only part of the story. For the baseline trained on clean text with seed 42, a net loss of 69 correct predictions at $20\%$ swaps is the residual of 142 newly induced errors offset by 73 previously misclassified articles that became correct after a swap suppressed a competing cue. Across topics, neural errors concentrate heavily in \emph{Business} (Table~\ref{tab:topic-correct-change}), which accounts for roughly $90\%$ of the net accuracy loss in both CNNs fitted with seed 42 as corrupted corporate or financial terms shift predictions toward \emph{Sci/Tech} and \emph{World}. Simultaneously, both CNNs classify more true \emph{World} articles correctly at $20\%$ swaps ($+28$ and $+31$ articles) while suffering drops in World $F_1$ of $2.98$ and $3.17$ percentage points, respectively, because 184 additional articles from other topics are misclassified as World in Subword CNN. In a practical news-routing pipeline, this asymmetric shift would flood the World feed with misrouted Business and Sci/Tech stories even as World recall appears to improve.

\subsection{Effectiveness of static swap augmentation}
\emph{Does static swap augmentation improve accuracy on corrupted articles while preserving accuracy on clean text?}

Replacing half of the training articles in each class with versions generated at the $10\%$ swap target yields a modest, positive mean gain for \textbf{Subword CNN} at the $20\%$ swap target: $+0.46$ percentage points on the original draw (pointwise 95\% bootstrap interval over articles $[+0.03, +0.90]$ percentage points) and $+0.39$ and $+0.31$ percentage points on the two additional draws (Table~\ref{tab:augmentation-draws}). By contrast, \textbf{Word TF-IDF LR} gains only $+0.08$ to $+0.14$ percentage points (which also reflects the validation-selected shift in regularization from $C = 2.0$ to $C = 8.0$), and \textbf{Word CNN} shows differences close to zero that change sign across draws ($-0.01$, $+0.12$, and $-0.04$ percentage points). Because Subword CNN composes token vectors from shared character fragments, training on corrupted words updates prefix, stem, and suffix buckets that generalize to unseen test misspellings even when the full word falls back to \texttt{<unk>}. Word CNN, by contrast, can benefit from a training typo only when that exact misspelled string enters the top 29,998 vocabulary entries and recurs verbatim at test time; otherwise, the token collapses to the single shared \texttt{<unk>} vector.

Nevertheless, two findings qualify the practical value of static replacement augmentation:

\paragraph{Variation across fitting seeds.}
Subword CNN's positive three-seed mean masks substantial run-to-run variation: seed 42 gained roughly $+1.0$ percentage point across all three $20\%$ swap draws, seed 43 gained $+0.36$ to $+0.37$ points in two draws and tied in the third, and seed 44 lost $0.03$ to $0.33$ points in every draw. Because each fitting seed varies both neural initialization/batch order and which 10,000 training articles are selected for corruption, static augmentation is not guaranteed to help on every individual training run.

\paragraph{Clean-text accuracy tradeoff.}
Replacing half of the clean training articles in place reduces how often the model observes canonical spellings during fitting. Although the three-seed augmentation differences on clean official test text have confidence intervals spanning zero, the separate 4,000-article confirmation check for Subword CNN (seed 42) shows that a $+0.60$ percentage point gain at $10\%$ swaps came alongside a $0.45$ percentage point loss on clean text, computed before rounding ($88.12\%$ to $87.68\%$, Table~\ref{tab:held-out-review}).

\subsection{Accuracy versus computational cost}
From an engineering standpoint, neither neural architecture justifies its substantially higher computational cost under this task and label budget. As shown in Table~\ref{tab:fit-cost-review}, \textbf{Word TF-IDF LR} fits its vectorizer and evaluates all three regularization candidates in a median of 13 seconds on four CPU threads, whereas \textbf{Word CNN} requires a median of 219 seconds (about $17\times$ as long) and \textbf{Subword CNN} requires a median of 642 seconds (about $50\times$ as long, based on unrounded medians) while still trailing the linear baseline by more than 2.7 percentage points on clean text and more than 3.4 percentage points at $20\%$ swaps. Within the neural family, Subword CNN with augmented training improves accuracy on corrupted text relative to the 8-epoch Word CNN, but at nearly triple the training time and still well behind the 13-second linear baseline.

\section{Schedule confounding, validity, and limitations}
A central methodological finding of our study is that the $1.30$ percentage point advantage of augmented Subword CNN over the original augmented Word CNN at $20\%$ swaps combines two distinct modifications: adding the 786,528-parameter character fragment table and extending the training schedule from 8 epochs (patience 2, fixed learning rate) to 16 epochs (patience 4, plateau learning rate halving). Our seed-42 schedule ablation (\textbf{Scheduled Word CNN}, Table~\ref{tab:scheduled-review}) isolates the optimization change while holding the 2,954,692-parameter Word CNN architecture fixed. On the official test set at $20\%$ swaps, granting Word CNN the 16-epoch schedule closes nearly the entire augmented gap ($85.50\%$ for Scheduled Word CNN vs.\ $85.57\%$ for Subword CNN and $84.75\%$ for the 8-epoch Word CNN) and slightly surpasses Subword CNN under clean training ($84.64\%$ vs.\ $84.49\%$).

Yet the validation learning curves in Figure~\ref{fig:learning} show that schedule length alone does not explain all differences between the two models: throughout epochs 2 to 8 in both training arms, prior to any learning rate reduction in either network, Subword CNN outperformed Scheduled Word CNN on the validation objective by up to $+1.78$ percentage points under identical learning rates and epoch counts. How much of Subword CNN's behavior reflects genuine subword generalization across misspelled stems (as illustrated in Table~\ref{tab:observed-swap}) versus extra embedding capacity or altered optimization dynamics therefore remains an open question requiring a matched-budget factorial experiment.

Several broader limitations also bound the scope of our conclusions:

\paragraph{Exploratory test exposure and selection bias.}
Because our initial baseline and Word CNN runs evaluated the official test split prior to designing the CNN validation screen, the subsequent neural comparisons should be interpreted as exploratory \citep{cawley2010}. Reusing a finite validation split across design iterations can likewise introduce mild selection optimism. We mitigated this risk by freezing all recipes prior to the final runs and evaluating the frozen seed-42 models once on a separate confirmation sample of 4,000 previously untouched articles, which reproduced the exact ranking of all three pipelines at $0\%$ and $10\%$ swaps.

\paragraph{Sources of evaluation uncertainty.}
Our pointwise 95\% intervals from the paired bootstrap over articles quantify test-set sampling variation across the 7,600 official test articles conditional on the fixed model checkpoints and swap draw. They do not capture variance across fresh training splits, larger seed ensembles (restricted to three fitting seeds for CPU tractability), or alternative noise models. Pairing bootstrap intervals with per-seed markers (Figure~\ref{fig:augmentation-review}) and two extra random swap draws (Table~\ref{tab:augmentation-draws}) provides a fuller view of run-to-run and corruption-draw variability.

\paragraph{Scope of the perturbation model.}
Transposing unequal adjacent interior letters preserves word length, word boundaries, and the first and last letters of every word. Real-world keyboard and OCR noise also includes character insertions, deletions, adjacent-key substitutions, and boundary errors that would disrupt prefix and suffix fragments (\texttt{<co}, \texttt{mer>}) left intact by interior swaps. Furthermore, because our nominal swap target $p$ applies only to eligible words ($\ge 4$ letters with unequal interior pairs), applying the same $p$ to corpora with shorter or longer words would alter a different share of total tokens; reporting realized token corruption ($12.93\%$ at $p = 20\%$) is therefore essential for cross-corpus comparison.

\paragraph{Single-label ambiguity and event overlap.}
AG News assigns a single topic label per article even when a story genuinely straddles multiple domains (such as telecommunications regulation or corporate sports sponsorships). Moreover, although our filter removed exact normalized text duplicates between the training and test splits, distinct wire stories covering the same underlying news event remain across splits; testing generalization to novel future events would require chronological or event-grouped partitioning.

\Needspace{8\baselineskip}
\section{Priorities for future work}
Building on these findings, we identify three focused priorities for follow-up work:

\paragraph{Matched-budget CNN factorial design.}
To cleanly separate character fragment representations from training schedule and parameter count, a confirmatory study should run the complete $2 \times 2$ factorial design in Table~\ref{tab:proposed-cnn-design} across all three fitting seeds and both training arms, adding the missing short-schedule Subword CNN control cell and equalizing total embedding parameters (for example, by adjusting word vocabulary size or embedding dimension so both architectures contain $\approx 3.0\mathrm{M}$ parameters).

\begin{table}[!htbp]
\centering\small
\begin{tabularx}{\linewidth}{@{}l>{\raggedright\arraybackslash}X>{\raggedright\arraybackslash}X@{}}
\toprule
Token representation & Short schedule (8 epochs, patience 2, constant LR) & Long schedule (16 epochs, patience 4, plateau LR decay) \\
\midrule
Words only & \textbf{Word CNN} (evaluated across 3 seeds) & \textbf{Scheduled Word CNN} (evaluated at seed 42; extend to 3 seeds) \\
Words + character fragments & \textbf{Short-schedule Subword CNN} (missing control cell) & \textbf{Subword CNN} (evaluated across 3 seeds) \\
\bottomrule
\end{tabularx}
\caption[Proposed factorial CNN comparison]{A $2 \times 2$ factorial design crossing token representation with training schedule. Running all four cells across matched seeds and equalized parameter budgets would isolate representational gains from optimization effects.}
\label{tab:proposed-cnn-design}
\end{table}

\paragraph{Targeted single-word perturbation probes.}
To directly test why global max pooling is more brittle than additive TF-IDF scoring, targeted single-word perturbation probes could swap one word at a time on fixed test articles, comparing logit-margin shifts when corrupting a word inside a CNN filter's winning pooling window (or holding a top TF-IDF weight) versus corrupting a non-winning word in the same document.

\paragraph{Dynamic epoch-level augmentation.}
Because static augmentation replaces half of the training corpus before epoch 1 (modifying 9,082 of the 10,000 selected articles at seed 42), edited training articles are never observed in their clean form. Comparing static replacement against dynamic per-epoch corruption, while varying the training corruption seed independently of the neural initialization seed, would test whether exposing every article in both clean and corrupted forms eliminates the clean-text accuracy tradeoff.

